\section{Principes fondamentaux}

\subsection{Espace des états}

\begin{omed}{Principe 1}
    À chaque système physique est associé un espace de Hilbert complexe $\mathcal{E}_H$, l’état physique du système étant entièrement défini par un ket normé $\ket{\psi} \in \mathcal{E}_H$.
\end{omed}

    On s’intéresse dans ce cours uniquement à des espaces de Hilbert dits \textit{séparables}, ce qui signifie qu’il admet au moins un ensemble dénombrable de kets $\ket{\psi_n}$ constituant une base orthonormée de l’espace.  

\subsection{Mesure}

On se place premièrement dans le cas \textbf{non-dégénéré}.

\begin{omed}{Principe 2}
    \begin{enumerate}[label=($p_{\arabic*}$)]
        \item À toute grandeur physique $A$ on associe un opérateur linéaire hermitien $\hat{A}$ de $\mathcal{E}_H$ appelé \textbf{observable}. On suppose ici que le spectre de $\hat{A}$ est \textit{discret} et \textit{non-dégénéré}.
        \item Une mesure de $A$ ne peut donner comme résultat qu’une valeur propre $a_n$ de $\hat{A}$.
        \item Pour un système se trouvant dans l’état $\ket{\psi}$, la probabilité de mesurer $a_n$ est
              \[ \mathbb{P}(a_n) := \abs{\spr{\psi_n}{\psi}}^2 \]
              où $\ket{\psi_n}$ est le vecteur propre associé à la valeur propre $a_n$.
        \item Si la mesure de $A$ donne $a_n$, alors le système est dans l’état $\ket{\psi_n}$ après mesure.
    \end{enumerate}
\end{omed}

Dans le cas général,

\begin{omed}{Principe 2}
    \begin{enumerate}[label=($p_{\arabic*}$)]
        \item À toute grandeur physique $A$ on associe un opérateur linéaire hermitien $\hat{A}$ de $\mathcal{E}_H$ appelé \textbf{observable}. On suppose ici que le spectre de $\hat{A}$ est de spectre discret.
        \item Une mesure de $A$ ne peut donner comme résultat qu’une valeur propre $a_n$ de $\hat{A}$.
        \item Pour un système se trouvant dans l’état $\ket{\psi}$, la probabilité de mesurer $a_n$ est
              \[ \mathbb{P}(a_n) := \bra{\psi} \hat{P}_n \ket{\psi} = \nor{\hat{P}_n \ket{\psi}}^2 \]
              où $\hat{P}_n$ est le projecteur sur le sous-espace propre associé à la valeur propre $a_n$.
        \item Si la mesure de $A$ donne $a_n$, alors le système après mesure est dans l’état
              \[ \ket{\psi'} = \frac{\hat{P}_n \ket{\psi}}{\nor{\hat{P}_n \ket{\psi}}} \]
    \end{enumerate}
\end{omed}

\subsection{Évolution temporelle}

\begin{omed}{Principe 3}
    L’évolution de l’état $\ket{\psi(t)}$ du système est gouvernée par l’\textbf{équation de Schrödinger}
    \[ i\hbar \frac{\text{d}\ket{\psi(t)}}{\text{d}t} = \hat{H}(t) \ket{\psi(t)} \]
    où $\bar{H}(t)$ est un observable appelé \textbf{Hamiltonien} et représentant l’énergie du système.
\end{omed}

On peut écrire l’état du système $\ket{\psi(t)}$ à l’instant $t$ à partir de l’état $\ket{\psi(t_0)}$ à l’instant $t_0$ à l’aide d’un opérateur unitaire $\hat{U}(t,t_0)$ appelé opérateur d’évolution selon l’expression 
\[ \ket{\psi(t)} = \hat{U}(t,t_0) \ket{\psi(t_0)} \]   
L’opérateur d’évolution obéit alors à l’équation différentielle 
\[ i\hbar \frac{\partial U(t,t_0)}{\partial t} = \hat{H}(t) \hat{U}(t,t_0) \]   
Dans le cas d’un hamiltonien indépendant du temps, la solution de cette équation s’écrit simplement 
\[ \hat{U}(t,t_0) = \hat{U}(t-t_0) = \exp\left(-i\frac{\hat{H}(t-t_0)}{\hbar}\right) \]  
ce qui donne la solution
\[ \ket{\psi(t)} = \exp\left(-i\frac{\hat{H}}{\hbar}(t - t_0)\right) \ket{\psi(t_0)} \]